% foundations/core_equations.tex

\chapter{Core Equations}
\label{chap:core-equations}

The equations below are executable knowledge: each relation names the
inputs it receives, the output it produces, the assumptions under which
that output is meaningful, and the questions it cannot decide. The
sequence follows the recurring CAD-curve and survey-polyline case from
construction, through realization and representation, to comparison.

% ============================================================
\section{Intrinsic State and Controls}
% ============================================================

The foundational spatial railway state is written as
\begin{equation}
x(s)
=
\bigl(\gamma(s),\theta(s),\phi(s)\bigr),
\label{eq:canonical-state}
\end{equation}
where \(\gamma(s)\) denotes planar position, \(\theta(s)\) heading, and
\(\phi(s)\) cant rotation.

The associated control pair is
\begin{equation}
u(s)
=
\bigl(\kappa(s),\omega(s)\bigr),
\qquad
\omega(s)=\phi'(s).
\label{eq:control-variables}
\end{equation}

The distinction between state and control follows the state model of
\cref{chap:state-model}: the state records the condition of the
alignment, while the controls govern its evolution along arc-length.

\paragraph{Execution contract.}
The equations receive a longitudinal domain, an initial state, and
admissible controls. They produce a state trajectory. They assume the
chosen state model and regularity conditions apply. They do not identify
the engineering object or determine the authority of any data source.

% ============================================================
\section{Planar Intrinsic Geometry}
% ============================================================

Arc-length parameterization gives
\begin{equation}
\gamma'(s)=\mathbf t(s),
\label{eq:state-evolution}
\end{equation}
and directional evolution is governed by
\begin{equation}
\theta'(s)=\kappa(s).
\label{eq:curvature-relation}
\end{equation}

In the planar case, the unit tangent is
\begin{equation}
\mathbf t(s)
=
\begin{pmatrix}
\cos\theta(s)\\
\sin\theta(s)
\end{pmatrix}.
\label{eq:planar-tangent}
\end{equation}

The canonical planar reconstruction system is therefore
\begin{equation}
\begin{aligned}
\gamma'(s)
&=
\begin{pmatrix}
\cos\theta(s)\\
\sin\theta(s)
\end{pmatrix},\\
\theta'(s)&=\kappa(s).
\end{aligned}
\label{eq:pose2-system}
\end{equation}

Given an initial position and heading, the curvature function determines
the planar alignment trajectory.

\paragraph{Execution contract.}
The reconstruction receives \(\kappa(s)\), \(\gamma(s_0)\), and
\(\theta(s_0)\); it produces \(\theta(s)\), \(\mathbf t(s)\), and
\(\gamma(s)\). It assumes arc-length parameterization and sufficient
regularity. A CAD curve sampled from \(\gamma\) is therefore a derived
representation. The reconstruction cannot decide whether a surveyed
polyline observes the same alignment or whether any discrepancy is
acceptable.

% ============================================================
\section{Cant Evolution and Foundational pose3 System}
% ============================================================

Cant rotation evolves according to
\begin{equation}
\phi'(s)=\omega(s).
\label{eq:cant-evolution}
\end{equation}

Together with the planar reconstruction equations, this gives the
foundational pose3 system
\begin{equation}
\begin{aligned}
\gamma'(s)
&=
\begin{pmatrix}
\cos\theta(s)\\
\sin\theta(s)
\end{pmatrix},\\
\theta'(s)&=\kappa(s),\\
\phi'(s)&=\omega(s).
\end{aligned}
\label{eq:pose3-system}
\end{equation}

Abstractly, the state evolution may be written as
\begin{equation}
x'(s)=f\bigl(x(s),u(s)\bigr).
\label{eq:general-state-system}
\end{equation}

This equation expresses the state-system form only. Profile, metric
realization, spatial frame construction, and physical interpretation are
introduced separately.

\paragraph{Execution contract.}
The system receives curvature and cant-rate controls plus an initial
state; it produces a pose trajectory. It assumes that the selected
pose model is applicable. It does not produce a physical realization or
an observation of one.

% ============================================================
\section{Railway Dynamic Relations}
% ============================================================

For the classical local engineering approximation, the effective
lateral acceleration is
\begin{equation}
a_{\mathrm{eff}}
=
v^2\kappa-g\sin\phi.
\label{eq:effective-lateral-acceleration}
\end{equation}

Ideal cant balance satisfies
\begin{equation}
v^2\kappa=g\sin\phi.
\label{eq:cant-equilibrium}
\end{equation}

For constant speed, differentiation with respect to time yields the
corresponding jerk relation
\begin{equation}
j
=
v^3\kappa'
-gv\cos\phi\,\phi'.
\label{eq:jerk-relation}
\end{equation}

These equations provide the foundational coupling between curvature,
cant, and operational velocity. Their railway interpretation and limits
are developed in the state and dynamics chapters.

\paragraph{Execution contract.}
These relations receive realized metric quantities and an applicable
local engineering approximation; they produce acceleration, equilibrium,
or jerk quantities. They assume the stated speed and differentiation
conditions. They evaluate physical consequences but cannot determine
operational admissibility or authorize a decision.

% ============================================================
\section{Track--World Relations}
% ============================================================

In a local planar realization, a track-related point with longitudinal
coordinate \(s\) and lateral offset \(d\) is mapped to world space by
\begin{equation}
x_{\mathrm w}(s,d)
=
\gamma(s)+d\,\mathbf n(s).
\label{eq:track-to-world}
\end{equation}

\begin{definition}[Track-to-world mapping]
\label{def:track-to-world}
The mapping defined by \cref{eq:track-to-world} relates track-related
coordinates to a world-space realization.
\end{definition}

The corresponding inverse problem is expressed as
\begin{equation}
(s^\ast,d^\ast)
=
\operatorname*{arg\,min}_{s,d}
\left\|
x_{\mathrm w}-\bigl(\gamma(s)+d\,\mathbf n(s)\bigr)
\right\|.
\label{eq:world-to-track}
\end{equation}

\begin{definition}[World-to-track mapping]
\label{def:world-to-track}
The world-to-track mapping determines an intrinsic track-related
interpretation of a world-space point within an applicable realization
context.
\end{definition}

A transformation between two world-space coordinate descriptions is
written abstractly as
\begin{equation}
x_2=\mathcal C(x_1).
\label{eq:coordinate-transformation}
\end{equation}

Coordinate transformation and world-to-track mapping are distinct
operations: the former changes world-space representation, while the
latter relates world space to an intrinsic track-related domain.

\paragraph{Execution contract.}
Track-to-world receives a realized alignment frame and track-related
coordinates; it produces world coordinates. World-to-track receives a
world point and the same applicable realization; it produces a candidate
intrinsic interpretation. Uniqueness requires a suitable local domain.
Neither operation establishes that two sources refer to the same object.

% ============================================================
\section{Representation and Engineering Realization}
% ============================================================

Sampling a continuous function at positions \(s_i\) is written as
\begin{equation}
\mathcal S(f)=\{f(s_i)\},
\label{eq:sampling}
\end{equation}
and reconstruction from sampled values as
\begin{equation}
\mathcal I\bigl(\{f(s_i)\}\bigr)=\widetilde f(s).
\label{eq:interpolation}
\end{equation}

An engineering realization operator relates an intended state to a
physically realized state:
\begin{equation}
\mathcal R:
\mathrm{pose3}_{\mathrm{target}}
\longrightarrow
\mathrm{pose3}_{\mathrm{real}}.
\label{eq:realization-operator}
\end{equation}

A schematic decomposition used in later realization chapters is
\begin{equation}
\mathcal R
=
\mathcal M\circ\mathcal I\circ\mathcal S,
\label{eq:core-realization-decomposition}
\end{equation}
where sampling, reconstruction, and physical response remain distinct
operator responsibilities.

At curvature level, the corresponding relation is
\begin{equation}
\mathcal R\bigl(\kappa_{\mathrm{target}}\bigr)
=
\kappa_{\mathrm{real}}.
\label{eq:curvature-realization}
\end{equation}

\begin{definition}[Curvature realization]
\label{def:curvature-realization}
Curvature realization denotes the relation between intended and
physically realized curvature under an applicable engineering
realization process.
\end{definition}

The associated inverse problem seeks an intended curvature satisfying
\begin{equation}
\mathcal R\bigl(\kappa_{\mathrm{target}}\bigr)
\approx
\kappa_{\mathrm{desired}}.
\label{eq:inverse-realization}
\end{equation}

\paragraph{Execution contract.}
Sampling receives a continuous description and sample locations;
interpolation receives sampled values and an interpolation rule;
realization receives a target condition and a physical-response model.
Their outputs retain those dependencies. In the test case, a survey
polyline can support an observed description of realized geometry, but
neither sampling nor interpolation turns the observation into physical
reality or gives it constructive authority.

% ============================================================
\section{Surface Curvature}
% ============================================================

For a curve constrained to a surface, total curvature decomposes into
geodesic and normal components:
\begin{equation}
\kappa^2=\kappa_g^2+\kappa_n^2.
\label{eq:curvature-decomposition}
\end{equation}

The geodesic curvature is expressed intrinsically by
\begin{equation}
\kappa_g
=
\left\|\nabla_s\mathbf t\right\|.
\label{eq:geodesic-curvature}
\end{equation}

The applicable surface, metric, connection, and engineering
interpretation are established in the ellipsoid and spatial-realization
chapters.

\paragraph{Execution contract.}
These equations receive a curve, surface, metric, and connection; they
produce curvature components in that realization. Without those inputs
the symbols are not executable. The components cannot decide which
surface model is authoritative for a project.

% ============================================================
\section{Optimization Relations}
% ============================================================

A schematic realization-aware optimization problem is
\begin{equation}
\min_u\;J\bigl(\mathcal R(x(u))\bigr).
\label{eq:general-optimization}
\end{equation}

For nonlinear least-squares problems, the objective takes the form
\begin{equation}
\min_x\frac12\|r(x)\|^2,
\label{eq:least-squares}
\end{equation}
with the Gauss--Newton approximation
\begin{equation}
H\approx J^T J.
\label{eq:gauss-newton}
\end{equation}

These equations specify common numerical structures only. Engineering
observations, residuals, constraints, admissibility, and decisions retain
their solver-independent meaning.

\paragraph{Execution contract.}
Optimization receives a declared objective, variables, residuals,
constraints, and any realization operator included in the model; it
produces a numerical candidate. Its validity depends on those choices
and on solver assumptions. A minimum is an evaluation result, not an
engineering decision.

% ============================================================
\section{Canonical Role}
% ============================================================

The equations now provide a dependency chain rather than a formula
catalogue. They can reconstruct a design trajectory, realize it in a
metric context, interpret observations, and compute qualified
differences. The next chapters provide the state and operator machinery
needed to execute that chain consistently. Engineering Concepts then
name the responsibilities that computation must preserve.
